% §1.1 Three Faces of a Noisy Bit
%
% Cross-reference resolution status:
%
% ALREADY DEFINED (within chapter 1, by this section):
%   \label{sec:three-faces}                 -- this section
%
% FORWARD REFS to later chapter 1 sections:
%   \Cref{sec:quantifying-noise}            -- §1.2
%   \Cref{sec:boolean-universe-first-look}  -- §1.3
%   \Cref{sec:sampling-rates}               -- §1.4
%   \Cref{sec:why-atom}                     -- §1.5
%
% FORWARD REFS to chapters 2-4:
%   \Cref{sec:bsc}                          -- chapter 2 §2.1
%   \Cref{sec:asymmetric-channel}           -- chapter 2 §2.2
%   \Cref{ch:axioms-kronecker}              -- chapter 4
%
% Monograph register: the three running threads (biased coin, BSC,
% classifier verdict) are introduced and the shared object named. They
% return at named points in chapters 2-4.

\section{Three Faces of a Noisy Bit}
\label{sec:three-faces}

Three systems from different corners of applied probability each produce a
single Boolean output that may be wrong, and at the level of input-output
behavior they are one object. We introduce them as running examples and
extract that object.

\subsection*{The biased coin}

Consider a coin whose bias depends on how it is held: palmed ``in'' it lands
heads with probability $\tau$, palmed ``out'' with probability $\epsilon$,
where $\tau$ and $\epsilon$ are neither equal nor near $1/2$. Read ``in'' as a
positive condition and ``out'' as a negative one. The coin reports each with
its own reliability: two inputs, one Boolean output, two independent
parameters. It returns as an asymmetric channel in
\Cref{sec:asymmetric-channel} and as the worked instance of the formal
definition in \Cref{ch:axioms-kronecker}.

\subsection*{The binary symmetric channel}

A single bit crosses a noisy line, a wire, a wireless link, a chip-to-chip
trace. For each bit independently the line flips the symbol with crossover
probability $\fprate$ and passes it through otherwise, with the same $\fprate$
in both directions: a $\top$ flips to $\bot$ and a $\bot$ to $\top$ with equal
probability. This is the binary symmetric channel (BSC), the simplest model of
a noisy one-bit line and the most studied channel in information theory; the
term \emph{channel} is Shannon's, for any single-input, single-output
stochastic transmission. The symmetry is a real constraint, since many
physical channels are asymmetric and need a separate rate per direction
(\Cref{sec:asymmetric-channel}). Chapter~2 gives the BSC a formal definition
and a matrix representation (\Cref{sec:bsc}).

\subsection*{The classifier verdict}

A spam filter emits a verdict on each message, spam ($\top$) or ham ($\bot$).
It flags ham as spam at the false-positive rate $\fprate$ and misses spam at
the false-negative rate $\fnrate$, and these are generally unequal: a filter
that drops legitimate mail at $1\%$ is a different product from one that leaks
$1\%$ of spam. Here the two parameters are forced by the problem rather than by
happenstance, since the two errors are distinct harms with distinct costs. A
score-thresholded classifier traces a curve of $(\fprate, \fnrate)$ pairs as
its threshold sweeps, the trade-off the ROC curve re-axes; no threshold drives
both rates to zero at once.

\subsection*{The shared object}

A coin held in or out, a wire carrying one bit, a filter emitting one verdict:
physical, informational, and algorithmic, yet identical in input-output
behavior. Each maps a Boolean input to a Boolean output that is correct with
some probability and wrong otherwise, and the error behavior is fixed by at
most two numbers, the rate of reporting $\top$ on a $\bot$-input and the rate
of reporting $\bot$ on a $\top$-input. The BSC carries one parameter because
its two rates are constrained equal; the coin and the filter carry two. Two
numbers describe everything probabilistic at the single-query level. We call
this object a \emph{noisy bit}.

The three faces used three vocabularies because each is natural in its own
letters, and they reconcile directly: the coin's $\tau$ is the true-positive
rate $\tprate = 1 - \fnrate$ and its $\epsilon$ is $\fprate$; the classifier's
$(\fprate, \fnrate)$ are the same quantities; the BSC is the special case
$\fprate = \fnrate$. Three commitments follow and hold throughout Part~I. The
noisy bit is \emph{atomic}: one Boolean input slot, one Boolean output slot,
nothing smaller, and everything larger (multi-bit channels, repeated-trial
algorithms such as Miller-Rabin, classifier ensembles) is built from noisy
bits, as chapters 2 and 4 show. It is \emph{behavioral}: the description fixes
the rates and says nothing about mechanism, so a biased coin, a thermal photon
detector, and a logistic-regression classifier with the same rates are the
same noisy bit. And it is \emph{parsimonious}: at most two numbers per noisy
bit. The rest of Part~I is the algebra of how these atoms compose.
